\section*{Prolog - Where Does What Disappears Go?}

My serious inquiry into existence and disappearance began one day when I was in my mid-twenties. At the time, I was a doctoral student. A younger friend whom I cared for deeply had recently graduated from university. Feeling that he was not yet fully prepared to seek employment, he asked to join our laboratory so that he could spend another year studying his field more seriously.

We had already known each other for several years through a student club. But the year we spent together in the laboratory brought us much closer. I gave him the seat beside mine and helped him organize his studies, consider his future, and prepare for employment. We saw each other almost every day, sharing trivial jokes, serious concerns, and the ordinary rhythms of laboratory life. During that year, he became much more than simply a younger colleague.

Eventually, he secured a position at one of the country’s leading companies. I was deeply pleased that the time we had spent preparing together had finally borne fruit. He left the laboratory and was just beginning the next stage of his life. Only a few months later, however, unexpected news arrived of an accident at the company.

He died.

That day, for the first time in my life, I watched someone close to me disappear inside a crematorium. When the cremation was over, his family members boarded a chartered bus and left the grounds one by one. I stood there, dazed, watching the bus depart. The place that had only moments earlier been filled with mourning voices and human movement was suddenly occupied by an immense emptiness.

White smoke was still rising from the crematorium chimney, slowly wavering as it climbed into the blue sky. Several birds I could not name were flying above it. As I watched them, a thought suddenly passed through my mind.

That smoke must contain part of him. The birds, too, were breathing. Part of him, now dispersed into the air as smoke, might enter their bodies with each breath. Something that had until only a short time ago belonged to a single human being might now be mingling with the bodies of several birds and flying across the blue sky.

At that moment, the scent of acacia blossoms unexpectedly reached me. I looked around. Acacia trees stood in a line along one side of a crumbling ocher-colored slope. Yet there were no flowers on them. Their branches were covered only with dense, dark-green leaves. Where, then, had the fragrance come from?

In my memory were the good moments I had shared with him. There was also the scent of acacia blossoms, a fragrance I had always loved. Perhaps those two memories had overlapped for an instant and returned to me as though they were a real sensation. Perhaps he, too, would remain somewhere beyond the visible flow of time and return without warning, just as that fragrance had returned to me.

Then another question arose.

Where, within my body, were all the experiences I had shared with him? Was memory merely a record of the past stored somewhere inside the brain? Or had those experiences already become part of my body and my way of thinking, continuing to live through the judgments I made and the actions I took?

If remembering him changed the way I treated someone else, would he not continue to affect other people through me? If so, could I really say that a person’s existence ended completely at the moment the body disappeared?

As these thoughts unfolded, everything around me began to appear more vivid than before. The grass, the trees, the birds crossing the sky, the dogs passing nearby, and the people exchanging greetings all seemed different.

How much had that grass grown through the previous spring and summer? What had the grass originally been? Matter from the soil must have entered through its roots. Sunlight and air must have passed through its leaves and become part of the plant now standing before me.

Long before this place became a crematorium, when it was still a field or a mountainside, perhaps a rabbit had once run across the same ground. After its death, its body might have mixed with the soil, moved with water, and been absorbed into countless grasses and trees. Those plants might then have withered and returned to the earth. Some might have become part of other animals, while other parts might have risen into the sky again as unnamed gases. 

Where, then, was that rabbit now? Had it ceased to exist? Or was it merely no longer bound together in the single form we once called a rabbit, its former substance now dispersed through grass, trees, air, water, and other living bodies? That day, I felt for the first time what it might truly mean to say that the world goes round and round.

It did not mean that the same being returned in exactly the same form. It meant that the things once gathered into one existence dispersed, mixed, and became parts of other existences. A fixed form vanished, but the elements that had constituted it and the effects it had produced continued through other structures.

Could this be what reincarnation really meant? Not the migration of an intact soul from one body into another, but the recognition that no existence is ever completely separate from the world—that everything gathered into one existence eventually disperses and becomes part of other forms of existence. After that day, I began to doubt the boundaries of existence.

Where does a person begin, and where does that person end? Is the surface of the body the true boundary of the self? At what point do the air one breathes and the food one consumes become part of that person? Are the memories and effects a person leaves behind unrelated to that person’s existence? Is death the disappearance of existence itself, or is it the dissolution of the organization that had temporarily held an existence together?

Questions about the past followed. If past experiences continue to transform who I am in the present, has the past truly disappeared? Is time an independent current that carries everything irreversibly into the past? Or is what we call the past already incorporated into the structure of the present, where it continues to operate?

At times, I wondered where he might be and what he might be doing now if he had not spent that year beside me preparing for employment. He might have made different preparations, joined a different company, and followed an entirely different path. And if the accident had never happened, what conversations would we still be having today? How might his choices and thoughts have shaped my present and continued to alter my future?

There could be no answer to such questions. Yet one thing became increasingly clear: an event does not necessarily end after it has passed. Not only what occurred, but also the possibilities that the event foreclosed, can leave traces in our present judgments and actions. The time I had spent with him, his death, and the questions that followed had already become part of the structure of who I had become.

I began to ask the same questions about space. Is the distance between me and a bird, between me and the grass, between me and the world, truly an empty interval? Are we independent beings separated from one another? Or are we continuously becoming parts of one another through countless exchanges and influences that remain invisible to us?

At first, these were only indistinct thoughts that arose in the face of a young friend’s death. But the questions did not disappear. They became broader and deeper with time. Human beings and other forms of life, matter and energy, memory and experience, time and space, motion and transformation, society and civilization—all the problems that had once seemed separate gradually began to converge into a single question.

Is independent existence truly possible? Or is everything we call an existence merely a temporary boundary maintained by countless differences, exchanges, and forms of cohesion? The white smoke rising from the crematorium chimney, the unnamed birds flying above it, and the sudden fragrance from acacia trees on which no blossoms could be seen remained within me for many years.

That day, I watched one existence disappear. Yet at that very place, I first began to wonder whether anything ever disappears completely. That experience became the trigger for my sustained inquiry into DISTANCE.

\pagebreak